% !PRINTED 381
\weberpage{381}

%% [running head, not body text -- commented out 2026-08-07] Page furniture (running
%% head, folio, signature) is excluded from this edition by standing policy.
%% \noindent\S.\,123. \hfill \textsc{Reduced Numbers.} \hfill 381

%<<CONT
\weq{6}{\omega = \alpha - \dfrac{1}{\omega_1},}

so that $\omega_1 = \xi_1 + i\eta_1$ is also equivalent to $\omega$. We now treat the
number $\omega_1$ again in the same way as $\omega$, by setting, if $\omega_1 - \alpha_1$
is not yet reduced,

\weq{7}{\omega_1 = \alpha_1 - \dfrac{1}{\omega_2}}

and so on. If we consider the sequence of the quantities formed on the analogy of (5),
\[
r_1^2 = (\xi_1-\alpha_1)^2 + \eta_1^2, \quad r_2^2 = (\xi_2-\alpha_2)^2 + \eta_2^2 \ldots,
\]
then all that now matters is to show that after a finite number of steps
of this kind we come to an $r_\nu^2$ that is equal to or greater
than 1.

But by (6) we have
\[
r^2 = (\omega-\alpha)(\omega'-\alpha) = \dfrac{1}{\xi_1^2+\eta_1^2},
\]
and from this the comparison of the imaginary parts on the two
sides of (6) gives:
\[
\eta = r^2 \eta_1.
\]
In the same way there follows the sequence of equations

\weq{8}{\eta = r^2\eta_1, \quad \eta_1 = r_1^2\eta_2, \quad \eta_2 = r_2^2\eta_3 \ldots}

Now, however, $i\eta = y\sqrt{d}$, and hence by the formula (6), \S.\,122
\[
i\eta = \dfrac{\sqrt{D}}{2c},
\]
in which $c$ is an integer. But because equivalent numbers have the same
discriminant, it follows in the same way:
\[
i\eta_1 = \dfrac{\sqrt{D}}{2c_1}, \quad i\eta_2 = \dfrac{\sqrt{D}}{2c_2}, \ldots,
\]
in which $c, c_1, c_2 \ldots$ is a sequence of positive integers. From this
one obtains from (8):

\weq{9}{c_1 = r^2 c, \quad c_2 = r_1^2 c_1, \quad c_3 = r_2^2 c_2 \ldots}

But so long as the numbers $r^2, r_1^2 \ldots r_\nu^2$ are less than 1, it follows from this

\weq{10}{c > c_1 > c_2 \ldots > c_{\nu+1},}

and because there can be only a finite number of positive numbers that are
less than $c$, this sequence must terminate, and our Theorem 1.
is thereby proved.
